% !PRINTED 17
\weberpage{17}
%<<CONT
isomorph sind, $\mu$-$\nu$-stufig isomorph zu einander nennt. Bei
Weitem der wichtigste ist der einstufige Isomorphismus, den wir
daher als Isomorphismus schlechtweg bezeichnen, während ein
mehrstufiger Isomorphismus immer durch einen Zusatz kenntlich
gemacht werden soll $^{1)}$.

\begin{center}
\S.\,6.

\gesperrt{Die Compositionsreihe und der Satz von C. Jordan.}
\end{center}

Eine Gruppe $P$ heisst \gesperrt{einfach}, wenn sie ausser sich selbst
und der Einheit keinen Normaltheiler hat, \gesperrt{zusammengesetzt},
wenn noch andere Normaltheiler vorhanden sind.

Ein Normaltheiler, der keinen anderen Normaltheiler \gesperrt{über
sich hat,} der also nicht Theil eines anderen echten Normaltheilers
 von $P$ ist, heisst \gesperrt{ein grösster Normaltheiler}
von $P$ $^{2)}$. Ist $P_1$ ein grösster Normaltheiler von $P$, $P_2$ ein
grösster Normaltheiler von $P_1$, $P_3$ ein grösster Normaltheiler
von $P_2$ u. s. f., so heisst die Reihe von Gruppen

\weq{1}{P,\, P_1,\, P_2,\, P_3 \ldots 1,}

die nothwendig mit der aus dem Einheitselemente allein gebildeten
 Gruppe 1 endigt, eine \gesperrt{Compositionsreihe} der
Gruppe $P$.

Wir wollen die Grade der Gruppen (1) mit $n, n_1, n_2, n_3 \ldots 1$
bezeichnen, und die Quotienten $n : n_1$, $n_1 : n_2$, $n_2 : n_3 \ldots$, also
die Indices von $P_1$ in Bezug auf $P$, $P_2$ in Bezug auf $P_1$ u. s. f.
mit $\nu_1, \nu_2, \nu_3 \ldots$ Die Reihe der ganzen Zahlen

\weq{2}{\nu_1,\, \nu_2,\, \nu_3 \ldots}

nennen wir die \gesperrt{Indexreihe} der Gruppe $P$.

Die grössten Normaltheiler $P_1, P_2, P_3 \ldots$ können für eine
gegebene Gruppe $P$ im Allgemeinen auf mehrere verschiedene
Arten ausgewählt werden, und danach können auch die Gradzahlen
 $n_1, n_2, n_3 \ldots$ und die Indices $\nu_1, \nu_2, \nu_3 \ldots$ verschieden
%CONT>>

$^{1)}$ Vgl. C.\,\gesperrt{Jordan}, Traité des substitutions. \gesperrt{Netto}, Substitutionentheorie
. Die Bezeichnung \glqq{}$\mu$-stufig\grqq{} rührt von \gesperrt{Netto} her. Nach \gesperrt{Jordan}
heisst der einstufige Isomorphismus \glqq{}holoëdrischer\grqq{}, der mehrstufige
\glqq{}meroëdrischer\grqq{} Isomorphismus. Der einstufige Isomorphismus wird bisweilen
 auch als Aequivalenz bezeichnet.

$^{2)}$ Ausgezeichnete Maximal-Untergruppe nach anderer Bezeichnung.

%% [page furniture, not body text -- commented out 2026-08-10] Page furniture (running
%% head, folio, signature) is excluded from this edition by standing policy.
%% Weber, Algebra. II. \hfill 2
